\section*{Bab \ref{ch:b2:curves} --- Kurva}

\begin{solution}{exo:b2:curves:1}
$\gamma'(t) = (1 - \cos t,\ \sin t)$, jadi
\[
\norm{\gamma'(t)}^2 = (1 - \cos t)^2 + \sin^2 t
= 2 - 2\cos t = 4\sin^2\tfrac t2 ,
\]
dan $\norm{\gamma'(t)} = 2\sin\frac t2$ (yang taknegatif pada $[0,
2\pi]$). Karena itu
\[
L = \int_0^{2\pi} 2\sin\tfrac t2\,\dd t
= \Bigl[-4\cos\tfrac t2\Bigr]_0^{2\pi} = 8 :
\]
jadi satu lengkung sikloidnya berpanjang $8$ (untuk roda berjari-jari
$1$) --- yakni hasil terkenal milik Wren, tanpa $\pi$ yang terlihat.
\end{solution}

\begin{solution}{exo:b2:curves:2}
Dengan $x = a\cos t$ dan $y = b\sin t$: berlaku $x' = -a\sin t$, $y' = b\cos
t$, $x'' = -a\cos t$, $y'' = -b\sin t$, jadi menurut
\cref{prop:b2:curves:kappaformula}
\[
\kappa(t)
= \frac{x'y'' - y'x''}{(x'^2 + y'^2)^{3/2}}
= \frac{ab\sin^2 t + ab\cos^2 t}
       {(a^2\sin^2 t + b^2\cos^2 t)^{3/2}}
= \frac{ab}{(a^2\sin^2 t + b^2\cos^2 t)^{3/2}} .
\]
Penyebutnya minimum ketika $\sin t = 0$ (bernilai $b^3$, yakni di titik
$(\pm a, 0)$) dan maksimum ketika $\cos t = 0$ (bernilai $a^3$, di titik
$(0, \pm b)$), karena $a > b$. Karena itu $\kappa$ maksimum di
ujung sumbu mayornya, $\kappa_{\max} = a/b^2$, dan minimum di
ujung sumbu minornya, $\kappa_{\min} = b/a^2$: jadi elipsnya
membengkok paling tajam di ujung sumbu \omterm{def:b2:curves:length}{panjangnya}.
\end{solution}

\begin{solution}{exo:b2:curves:3}
Untuk grafik $\gamma(x) = (x, \cosh x)$: berlaku $\norm{\gamma'(x)} =
\sqrt{1 + \sinh^2 x} = \cosh x$, jadi \omterm{def:b2:curves:length}{panjang busurnya} dari $0$ sampai $x$
adalah $\int_0^x \cosh u\,\dd u = \sinh x$. Adapun \omterm{thm:b2:curves:frenet2d}{kelengkungan} sebuah grafik
(\cref{prop:b2:curves:kappaformula}):
\[
\kappa(x) = \frac{f''(x)}{(1 + f'(x)^2)^{3/2}}
= \frac{\cosh x}{\cosh^3 x} = \frac{1}{\cosh^2 x} ,
\]
jadi $R(x) = \cosh^2 x$, seperti diumumkan. Perhatikan kebetulan yang rapi
$R(x) = 1 + s(x)^2$ dengan $s = \sinh x$ sebagai \omterm{def:b2:curves:length}{panjang busurnya}: jadi
\omterm{def:b2:curves:curvature}{jari-jari kelengkungan} katenarinya tumbuh dengan kuadrat \omterm{def:b2:curves:length}{panjang
busurnya} dari puncaknya.
\end{solution}

\begin{solution}{exo:b2:curves:4}
$\gamma'(t) = e^t(\cos t - \sin t,\ \sin t + \cos t)$, jadi
\[
\langle \gamma(t), \gamma'(t)\rangle = e^{2t}
\bigl(\cos t(\cos t - \sin t) + \sin t(\sin t + \cos t)\bigr)
= e^{2t},
\]
sedangkan $\norm{\gamma(t)} = e^t$ dan $\norm{\gamma'(t)} =
e^t\sqrt 2$. Karena itu
\[
\cos\angle\bigl(\gamma, \gamma'\bigr)
= \frac{e^{2t}}{e^t \cdot e^t\sqrt2} = \frac{1}{\sqrt2} :
\]
jadi singgungnya selalu bersudut $\pi/4$ terhadap jari-jarinya --- yakni
sifat \emph{sudut-sama} milik spiral logaritmiknya. Adapun \omterm{def:b2:curves:length}{panjang busurnya}
pada $(-\infty, 0]$:
\[
\int_{-\infty}^0 \norm{\gamma'(t)}\,\dd t
= \sqrt2\int_{-\infty}^0 e^t\,\dd t = \sqrt 2 ,
\]
yang berhingga walaupun spiralnya melilit tak berhingga kali mengelilingi
titik asalnya. Adapun \omterm{thm:b2:curves:frenet2d}{kelengkungannya}: dengan $x'y'' - y'x''$ yang dihitung dari $\gamma'' =
e^t(-2\sin t,\ 2\cos t)$,
\[
x'y'' - y'x'' = e^{2t}\bigl(2\cos t(\cos t - \sin t)
+ 2\sin t(\sin t + \cos t)\bigr) = 2e^{2t},
\]
jadi $\kappa(t) = \dfrac{2e^{2t}}{(e^t\sqrt2)^3} =
\dfrac{1}{e^t\sqrt2}$: sehingga \omterm{thm:b2:curves:frenet2d}{kelengkungannya} $1/(\sqrt2\,
\norm{\gamma})$, yang meluruh ketika spiralnya tumbuh.
\end{solution}

\begin{solution}{exo:b2:curves:5}
\emph{Busur pertamanya:} $\gamma(t) = (t^2,\ t^4 + t^5)$. Turunannya di
$0$: $\gamma'' = (2, 0) \neq 0$, jadi $p = 2$. Lalu $\gamma^{(3)}(0)
= (0, 0)$ dan $\gamma^{(4)}(0) = (0, 24)$, yang tak segaris dengan $(2,
0)$: jadi $q = 4$. Keduanya genap: yakni \emph{titik runcing jenis kedua} --- sebab kedua
cabangnya pergi ke arah $+u = (1,0)$ lalu tinggal di sisi yang sama
dari singgungnya. (Memang $y = x^2 \pm x^{5/2}$ pada kedua
cabangnya: jadi bertanda sama untuk $x$ yang kecil.)

\emph{Busur keduanya:} $\gamma(t) = (t^3, t^4)$. Berlaku $\gamma'(0) =
\gamma''(0) = 0$ dan $\gamma^{(3)}(0) = (6, 0)$: jadi $p = 3$, yang ganjil. Berikutnya
$\gamma^{(4)}(0) = (0, 24)$: jadi $q = 4$, yang genap. Ganjil--genap:
yakni \emph{titik biasa} --- sebab kendati kecepatannya lenyap,
lintasan $y = x^{4/3}$ tetap menyeberangi titik asalnya dengan mulus, sambil tinggal
di atas singgungnya $y = 0$.
\end{solution}

\begin{solution}{exo:b2:curves:6}
Turunkan $c(s) = \gamma(s) + \dfrac{1}{\kappa(s)}N(s)$ dengan memakai
rumus Frenet bidangnya
(\cref{thm:b2:curves:frenet2d}):
\[
c'(s) = T(s) - \frac{\kappa'(s)}{\kappa(s)^2}N(s)
+ \frac{1}{\kappa(s)}\,\bigl(-\kappa(s)T(s)\bigr)
= -\frac{\kappa'(s)}{\kappa(s)^2}\,N(s) ,
\]
sebab suku singgungnya saling meniadakan tepat. Jadi kecepatan
evolutnya diusung $N(s)$, yang mengarahkan garis normal
$\gamma$ di $\gamma(s)$ --- dan titik $c(s)$ terletak pada
garis normal itu juga: jadi evolutnya adalah \emph{\omterm{pb:b2:curves:1}{selubung} normalnya}.
(Adapun di tempat $\kappa' = 0$ evolutnya punya titik singular; dan inilah
yang menghasilkan titik runcing pada evolut sebuah elips.)
\end{solution}

\begin{solution}{exo:b2:curves:7}
Misalkan $G(s) = F(s)F(s)^{\mathsf T}$. Maka, dengan memakai $F' = FA$ dan
$(F^{\mathsf T})' = (F')^{\mathsf T} = A^{\mathsf T}F^{\mathsf T}$,
\[
G' = F'F^{\mathsf T} + F(F^{\mathsf T})'
= FAF^{\mathsf T} + FA^{\mathsf T}F^{\mathsf T}
= F(A + A^{\mathsf T})F^{\mathsf T} = 0
\]
berkat keantisimetrikannya. Jadi $G$ tetap pada intervalnya, dan sama dengan
$G(s_0) = F(s_0)F(s_0)^{\mathsf T} = I$: sehingga $F(s)$ ortogonal untuk
setiap $s$. (Atau, tanpa menghitung $G'$ menjadi nol: sebab baik
$G$ maupun $I$ yang tetap menyelesaikan sistem linear $Y' = YA +
A^{\mathsf T}Y$ dengan nilai awal yang sama, lalu ketunggalan Cauchy--Lipschitz
bagi sistem linear, \cref{ch:b2:diffeq}, memaksa $G
\equiv I$.)

\emph{Kaitannya:} sistem Frenet $(T, N, B)' = (T, N, B)\,A(s)$
punya matriks koefisien yang antisimetrik
\[
A = \begin{pmatrix} 0 & -\kappa & 0\\ \kappa & 0 & -\tau\\
0 & \tau & 0\end{pmatrix}
\]
(dengan kolomnya menyatakan $T', N', B'$). Adapun perhitungan di atas menunjukkan
bahwa kerangka penyelesaian yang bermula ortonormal \emph{tetap}
ortonormal --- yakni langkah kunci pada teorema dasar yang
membangun ulang sebuah kurva dari $(\kappa, \tau)$.
\end{solution}

\begin{solution}{exo:b2:curves:8}
Menurut \cref{thm:b2:curves:fundamental} (yakni bagian ketunggalannya), ada
fungsi sudut $\mathcal{C}^1$ $\varphi$ dengan $T(s) =
(\cos\varphi(s), \sin\varphi(s))$ dan $\varphi' = \kappa$. Karena itu
\[
\int_0^L \kappa(s)\,\dd s = \varphi(L) - \varphi(0) .
\]
Karena busurnya tertutup berkala $L$, berlaku $T(L) = T(0)$:
jadi $(\cos\varphi(L), \sin\varphi(L)) = (\cos\varphi(0),
\sin\varphi(0))$, sehingga $\varphi(L) - \varphi(0) \in 2\pi\Z$. Jadi
\omterm{thm:b2:curves:frenet2d}{kelengkungan} total sebuah kurva tertutup selalu berupa kelipatan bulat
$2\pi$ --- dengan bilangan bulatnya berupa \emph{bilangan lilitan}
singgungnya (yakni banyaknya putaran penuh yang dibuat $T$). Untuk
lingkaran berjari-jari $R$: berlaku $\kappa = 1/R$ dan $L = 2\pi R$, jadi \omterm{thm:b2:curves:frenet2d}{kelengkungan}
totalnya $2\pi$ dengan bilangan lilitan $1$; dan teorema singgung
berputar menyatakan bahwa nilai ini berlaku bagi setiap kurva tertutup sederhana.
\end{solution}

\begin{solution}{exo:b2:curves:9}
Karena $\tau \equiv 0$, kurvanya terletak pada sebuah bidang
(\cref{prop:b2:curves:torsion}); jadi bekerjalah di bidang itu. Tinjaulah
calon pusatnya
\[
c(s) = \gamma(s) + \frac{1}{\kappa}N(s)
\qquad (\kappa \text{ tetap}).
\]
Menurunkannya dengan rumus Frenet ($N' = -\kappa T + \tau B
= -\kappa T$ di sini) memberi
\[
c'(s) = T + \frac1\kappa(-\kappa T) = 0 ,
\]
jadi $c$ merupakan titik tetap $\Omega$. Maka $\norm{\gamma(s) -
\Omega} = \norm{-\frac1\kappa N(s)} = \frac1\kappa$ untuk setiap $s$:
sehingga kurvanya terletak pada lingkaran berpusat $\Omega$ dan berjari-jari
$1/\kappa$ (di bidangnya), lalu karena ia busur yang tak tetap padanya, ia
merupakan busur lingkaran itu.
\end{solution}

\begin{solution}{exo:b2:curves:10}
Untuk grafik $\gamma(x) = (x, x^2/2)$, berlaku $\norm{\gamma'(x)} =
\sqrt{1 + x^2}$, jadi $L = \int_0^a\sqrt{1+x^2}\,\dd x$.
Setelah menyulihkan $x = \sinh u$ (dengan $\dd x = \cosh u\,\dd u$, dan $u$ dari
$0$ ke $u_a = \ln(a + \sqrt{1+a^2})$):
\[
L = \int_0^{u_a}\cosh^2 u\,\dd u
= \frac12\bigl[u + \sinh u\cosh u\bigr]_0^{u_a}
= \frac12\Bigl(\ln\bigl(a + \sqrt{1+a^2}\bigr) +
a\sqrt{1+a^2}\Bigr),
\]
dengan memakai $\cosh^2 u = \frac{1 + \cosh 2u}2$ beserta $\sinh u_a = a$ dan
$\cosh u_a = \sqrt{1 + a^2}$.
\end{solution}

\begin{solution}{exo:b2:curves:11}
Parameterkan lewat \omterm{def:b2:curves:length}{panjang busurnya}
(\cref{thm:b2:curves:arclength}) lalu tetapkan $\lambda(s) =
\langle P - \gamma(s), T(s)\rangle$, yakni fungsi $\mathcal C^1$;
lalu karena $P$ terletak pada \omterm{def:b2:curves:arc}{garis singgung} di $\gamma(s)$,
vektor $P - \gamma(s)$ segaris dengan $T(s)$, jadi $P =
\gamma(s) + \lambda(s)T(s)$. Setelah diturunkan,
\[
0 = T(s) + \lambda'(s)T(s) + \lambda(s)T'(s)
= \bigl(1 + \lambda'(s)\bigr)T(s) + \lambda(s)T'(s),
\]
dan $T'(s) \perp T(s)$ (turunkan $\norm T^2 = 1$), jadi
kedua komponennya lenyap: $\lambda' = -1$ dan $\lambda T' = 0$.
Lalu $\lambda(s) = c - s$ lenyap paling banyak sekali, jadi $T' = 0$
pada sebuah himpunan padat, sehingga di mana-mana berkat kekontinuannya: jadi $T$ merupakan
vektor satuan yang tetap dan $\gamma(s) = \gamma(s_0) + (s -
s_0)T$: yakni sebuah garis lurus (lewat $P$, sebagaimana mestinya).
\end{solution}

\begin{solution}{exo:b2:curves:12}
(a) Matriks Frenetnya
\[
A(s) = \begin{pmatrix} 0 & -\kappa & 0\\ \kappa & 0 & -\tau\\
0 & \tau & 0\end{pmatrix}
\]
punya entri yang \omterm{def:b2:metric:continuity}{kontinu}, jadi sistem linear $F' = FA(s)$ dengan
$F(s_0) = F_0$ (yakni matriks ortonormal yang langsung) punya satu penyelesaian
tunggal pada seluruh $J$ (\cref{thm:b2:diffeq:linear}). Lalu menurut
\cref{exo:b2:curves:7}, $F(s)$ ortogonal untuk setiap $s$;
sedangkan $\det F$ \omterm{def:b2:metric:continuity}{kontinu} bernilai di $\{\pm1\}$ dan sama dengan
$1$ di $s_0$, jadi $F(s)$ langsung untuk setiap $s$.

(b) Bacalah baris $T, N, B$ milik $F$ (sehingga $T' = \kappa
N$, $N' = -\kappa T + \tau B$, $B' = -\tau N$) lalu tetapkan
$\gamma(s) = \gamma_0 + \int_{s_0}^s T(u)\,\dd u$. Maka
$\gamma' = T$ merupakan vektor satuan: jadi berkelajuan satuan; lalu $T' = \kappa N$
dengan $\kappa > 0$ dan $N$ satuan yang ortogonal terhadap $T$, jadi $\gamma$
bersifat biregular dengan \omterm{thm:b2:curves:frenet2d}{kelengkungan} $\norm{T'} = \kappa$ dan
normal utama $N$; sedangkan binormalnya adalah $T \wedge N = B$
(yakni kerangka ortonormal yang langsung), dan $B' = -\tau N$ mengenali
\omterm{thm:b2:curves:frenet3d}{torsinya} sebagai $\tau$.

(c) Misalkan $\gamma_1, \gamma_2$ kurva biregular berkelajuan satuan
dengan $(\kappa, \tau)$ yang sama. Maka ada satu isometri langsung
$\Phi = \rho + w$ (dengan $\rho \in SO(3)$) yang tunggal, yang mengirim
$\gamma_1(s_0)$ ke $\gamma_2(s_0)$ dan kerangka Frenet
$\gamma_1$ di $s_0$ ke kerangka $\gamma_2$ di $s_0$. Adapun kurva
$\Phi\circ\gamma_1$ berkelajuan satuan dengan invarian yang sama
(sebab kerangkanya adalah $\rho$ yang diterapkan pada kerangka $\gamma_1$, dan
$\rho$ memelihara hasil kali silangnya karena ia langsung). Kini
kerangka $\Phi\circ\gamma_1$ dan $\gamma_2$ sama-sama menyelesaikan $F' =
FA(s)$ dengan nilai awal yang sama, jadi keduanya berimpit berkat
ketunggalannya; khususnya singgungnya sepakat, lalu setelah diintegralkan
dari titik bersamanya $s_0$: $\Phi\circ\gamma_1 = \gamma_2$.
\end{solution}

\begin{solution}{pb:b2:curves:1}
\textbf{1.} Sistemnya linear dalam $(x, y)$ dengan \omterm{def:b2:linalg:det}{determinan}
$\Delta(t) = a b' - a'b \neq 0$: jadi aturan Cramer memberi
penyelesaian tunggalnya
\[
x = \frac{c b' - c' b}{\Delta}, \qquad
y = \frac{a c' - a' c}{\Delta}.
\]

\textbf{2.} Karena $a(t)x(t) + b(t)y(t) = c(t)$ secara identik,
menurunkannya memberi $a'x + b'y + ax' + by' = c'$; lalu
persamaan karakteristik keduanya membunuh $a'x + b'y - c'$, jadi
$a(t)x'(t) + b(t)y'(t) = 0$: sehingga $E'(t)$ ortogonal terhadap $(a,
b)$, jadi sejajar dengan $(-b, a)$, yakni arah $D_t$. Lalu karena
$E(t) \in D_t$ (menurut persamaan pertamanya) dan $E'(t) \neq 0$, garis
$D_t$ persis merupakan \omterm{def:b2:curves:arc}{garis singgung} kurva $E$ di $E(t)$.

\textbf{3.} Di sini $(a, b, c) = (t, -1, t^2/2)$, jadi $\Delta =
t\cdot0 - 1\cdot(-1) = 1$ dan
\[
x = \frac{c b' - c' b}{\Delta} = \tfrac{t^2}2\cdot 0 +
t = t, \qquad
y = \frac{a c' - a' c}{\Delta} = t\cdot t - \tfrac{t^2}2 =
\tfrac{t^2}2 :
\]
jadi \omterm{pb:b2:curves:1}{selubung} \omterm{def:b2:curves:arc}{garis singgung} parabolanya adalah
parabolanya sendiri, sebagaimana mestinya.

\textbf{4.} Garis normalnya adalah $\langle M, T(s)\rangle =
\langle\gamma(s), T(s)\rangle$: dengan koefisien $a = T_1$, $b =
T_2$, $c = \langle\gamma, T\rangle$. Menurunkannya dengan
Frenet ($T' = \kappa N$): $a' = \kappa N_1$, $b' = \kappa
N_2$, dan $c' = \langle T, T\rangle + \langle\gamma, \kappa
N\rangle = 1 + \kappa\langle\gamma, N\rangle$. Lalu persamaan
karakteristik keduanya $\kappa\langle M, N\rangle = 1 +
\kappa\langle\gamma, N\rangle$ berbunyi $\kappa\langle M -
\gamma, N\rangle = 1$. Sedangkan yang pertama mengatakan $M - \gamma \perp T$,
jadi $M - \gamma = \mu N$ dengan $\mu = 1/\kappa$: sehingga
titik karakteristiknya adalah $\gamma + \frac1\kappa N$, yakni \omterm{def:b2:curves:curvature}{pusat
kelengkungannya}, dan \omterm{pb:b2:curves:1}{selubung} normalnya adalah evolut
pada \cref{exo:b2:curves:6}. Akhirnya $\Delta = T_1\kappa N_2 -
\kappa N_1 T_2 = \kappa\det(T, N) = \kappa \neq 0$.

\textbf{5.} Berkasnya: sistem $x\cos\theta + y\sin\theta =
0$ dan $-x\sin\theta + y\cos\theta = 0$ berdeterminan $1$ dengan
penyelesaian $(0,0)$ untuk setiap $\theta$: jadi $E \equiv (0,0)$, $E'
\equiv 0$, dan tak ada kurva --- melainkan sekadar titik bersama
semua garisnya. Adapun keluarga sejajarnya: $a' = b' = 0$ memberi
$\Delta \equiv 0$ dan persamaan keduanya $0 = c'(t)$, yang
gagal begitu keluarganya sungguh bergerak: jadi tak ada titik
karakteristiknya, dan memang keluarga garis sejajar tak menyentuh
sebuah kurva sepanjang semua anggotanya.

\textbf{6.} Garis lewat $(\cos t, 0)$ dan $(0, \sin t)$
adalah $\frac x{\cos t} + \frac y{\sin t} = 1$, yakni $x\sin t +
y\cos t = \sin t\cos t$. Dengan $(a, b, c) = (\sin t, \cos t,
\sin t\cos t)$: berlaku $a' = \cos t$, $b' = -\sin t$, $c' = \cos
2t$, dan $\Delta = -\sin^2 t - \cos^2 t = -1$. Lalu Cramer:
\begin{align*}
x &= \frac{cb' - c'b}{-1} = \sin^2 t\cos t + \cos 2t\cos t
= \cos t\,(\sin^2 t + \cos^2 t - \sin^2 t) = \cos^3 t,\\
y &= \frac{ac' - a'c}{-1} = \sin t\cos^2 t - \sin t\cos 2t
= \sin t\,(\cos^2 t - \cos^2 t + \sin^2 t) = \sin^3 t .
\end{align*}
Dan $(\cos^3t)^{2/3} + (\sin^3t)^{2/3} = 1$: yakni \omterm{pb:b2:curves:1}{astroidnya}.

\textbf{7.} Turunan $E'(t) = 3(-\cos^2 t\sin t,\ \sin^2 t\cos t)$
lenyap di $t = 0$. Di sana, $E''(0) = (-3, 0) \neq 0$ memberi
$p = 2$; lalu komponen-$x$ milik $E$ bersifat genap dalam $t$, jadi
$E'''(0) = (0, 6)$ yang tak segaris: $q = 3$. Genap--ganjil: jadi titik runcing
jenis pertama di $(1, 0)$
(\cref{prop:b2:curves:local}), yang bersinggung sepanjang
sumbu-$x$. Adapun kesimetrian $x \mapsto -x$, $y \mapsto -y$, dan $(x,
y) \mapsto (y, x)$ milik \omterm{pb:b2:curves:1}{astroidnya} mengangkut titik runcingnya ke $(-1,
0)$ dan $(0, \pm1)$.

\textbf{8.} Berlaku $E'(t) = 3\sin t\cos t\,(-\cos t, \sin t)$, jadi
$\norm{E'(t)} = 3\abs{\sin t\cos t} = \tfrac32\abs{\sin 2t}$.
Satu kuadrannya: $\int_0^{\pi/2}\tfrac32\sin 2t\,\dd t =
\tfrac32$, jadi berkat kesimetriannya \omterm{def:b2:curves:length}{panjang} totalnya adalah $4 \cdot
\tfrac32 = 6$.

\textbf{9.} Berlaku $E(t) - P_t = (\cos^3 t - \cos t,\ \sin^3 t) =
\sin^2 t\,(-\cos t,\ \sin t) = \sin^2 t\,(Q_t - P_t)$. Jadi
titik sentuhnya adalah \omterm{def:b2:affine:barycenter}{barisentrum} $(P_t, \cos^2 t)$ dan
$(Q_t, \sin^2 t)$: sehingga ketika $t$ berjalan dari $0$ ke $\pi/2$ ia meluncur
dari ujung lantai ke ujung dinding tangganya.

\textbf{10.} Pada kuadran pertamanya, daerah di bawah
\omterm{pb:b2:curves:1}{astroidnya} berluas $\int_0^1 y\,\dd x$ dengan $x = \cos^3 t$
yang turun dari $1$ ke $0$ ketika $t$ berjalan dari $0$ ke $\pi/2$:
\[
\int_0^1 y\,\dd x
= \int_{\pi/2}^{0}\sin^3 t\,(-3\cos^2 t\sin t)\,\dd t
= 3\int_0^{\pi/2}\sin^4 t\cos^2 t\,\dd t .
\]
Linearkan: $\sin^4 t\cos^2 t = (\sin t\cos t)^2\sin^2 t =
\tfrac18\bigl(\sin^2 2t - \sin^2 2t\cos 2t\bigr)$, lalu
$\int_0^{\pi/2}\sin^2 2t\,\dd t = \tfrac\pi4$ sedangkan
$\int_0^{\pi/2}\sin^2 2t\cos 2t\,\dd t =
\bigl[\tfrac{\sin^3 2t}6\bigr]_0^{\pi/2} = 0$. Jadi
integralnya $\tfrac\pi{32}$, luas kuadrannya
$\tfrac{3\pi}{32}$, dan luas yang dilingkupinya $4 \cdot
\tfrac{3\pi}{32} = \tfrac{3\pi}8$.

\textbf{11.} Singgung di $\gamma(t) = (t, t^2/2)$
diarahkan $(1, t)$, jadi garis normalnya adalah $\{(x, y) : (x -
t) + t\,(y - t^2/2) = 0\}$, yakni $x + t\,y = t +
\tfrac{t^3}2$.

\textbf{12.} Berlaku $(a, b, c) = (1, t, t + t^3/2)$: $a' = 0$, $b' =
1$, $c' = 1 + \tfrac32 t^2$, dan $\Delta = 1$. Maka $y = ac' -
a'c = 1 + \tfrac32t^2$ dan
\[
x = cb' - c'b = t + \tfrac{t^3}2 - t\Bigl(1 +
\tfrac32t^2\Bigr) = -t^3 .
\]
Setelah $t$ dihilangkan: $t^2 = \tfrac23(y - 1)$ dan $x^2 = t^6 =
\tfrac8{27}(y - 1)^3$: yakni sebuah parabola semikubik bertitik puncak
$(0, 1)$.

\textbf{13.} Berlaku $T = (1, t)/\sqrt{1+t^2}$, $N = (-t,
1)/\sqrt{1+t^2}$, dan $\kappa = (1+t^2)^{-3/2}$, jadi
\[
\gamma + \frac1\kappa N = (t, \tfrac{t^2}2) +
(1+t^2)\,(-t, 1) = \Bigl(-t^3,\ 1 + \tfrac32t^2\Bigr),
\]
yakni kurva yang sama: jadi \omterm{pb:b2:curves:1}{selubung} normalnya adalah tempat kedudukan
\omterm{def:b2:curves:curvature}{pusat kelengkungannya}, sebagaimana dijanjikan pertanyaan 4.

\textbf{14.} Turunan $E'(t) = (-3t^2, 3t)$ lenyap di $t = 0$;
lalu $E''(0) = (0, 3) \neq 0$ memberi $p = 2$ dan $E'''(0) = (-6,
0)$ memberi $q = 3$: jadi titik runcing jenis pertama di $(0, 1)$. Adapun
puncaknya adalah tempat $\kappa = (1+t^2)^{-3/2}$ maksimum, jadi
$\kappa'(0) = 0$ sehingga kecepatan evolutnya
$-\frac{\kappa'}{\kappa^2}N$ lenyap persis di sana: jadi titik runcing
evolutnya duduk pada ekstremum \omterm{thm:b2:curves:frenet2d}{kelengkungannya}.

\textbf{15.} Normal di parameter $t$ melalui
$(x_0, y_0)$ bila dan hanya bila $x_0 + t\,y_0 = t + \tfrac{t^3}2$, yakni
\[
\frac{t^3}2 + (1 - y_0)\,t - x_0 = 0 ,
\]
yakni sebuah kubik dalam $t$: jadi satu atau tiga akar real (yang dicacah tanpa
multiplisitas, bagi titik yang generik). Pada sumbunya $x_0 = 0$ ia
terfaktorkan sebagai $t\bigl(\tfrac{t^2}2 + 1 - y_0\bigr) = 0$: yakni
akar $t = 0$ (sebab sumbunya adalah normal di puncaknya), ditambah
$t = \pm\sqrt{2(y_0 - 1)}$ ketika $y_0 > 1$. Jadi: satu normal
bila $y_0 < 1$, tiga bila $y_0 > 1$, dan di $y_0 = 1$
akar rangkap tiganya menandai titik runcing evolutnya. Secara umum sebuah
akar ganda kubiknya berarti titiknya memenuhi baik
persamaan garisnya maupun turunan-$t$-nya --- jadi ia terletak \emph{pada
\omterm{pb:b2:curves:1}{selubungnya}}: sehingga evolutnya persis merupakan perbatasan antara
daerah satu-normal dan daerah tiga-normalnya.

\textbf{16.} Pemantulan pada cerminnya membalik komponen normal
arahnya dan mempertahankan yang tangensialnya:
dengan menulis $u = \langle u, n\rangle n + u_{\mathrm{tan}}$,
arah pantulnya adalah $u_{\mathrm{tan}} - \langle u,
n\rangle n = u - 2\langle u, n\rangle n$. Di sini $\langle u,
n\rangle = \cos\theta$, jadi
\[
v = (1, 0) - 2\cos\theta\,(\cos\theta, \sin\theta)
= (1 - 2\cos^2\theta,\ -2\sin\theta\cos\theta)
= -(\cos2\theta,\ \sin2\theta).
\]

\textbf{17.} Sinar pantulnya melalui $P_\theta =
(\cos\theta, \sin\theta)$ dengan arah $(\cos2\theta,
\sin2\theta)$; jadi sebuah vektor normalnya adalah $(-\sin2\theta,
\cos2\theta)$, sehingga garisnya adalah
\[
-\sin2\theta\,(x - \cos\theta) + \cos2\theta\,(y -
\sin\theta) = 0,
\]
dan konstantanya adalah $-\sin2\theta\cos\theta +
\cos2\theta\sin\theta = -\sin\theta$: jadi setelah dikalikan $-1$,
$x\sin2\theta - y\cos2\theta = \sin\theta$.

\textbf{18.} Berlaku $(a, b, c) = (\sin2\theta, -\cos2\theta,
\sin\theta)$: $a' = 2\cos2\theta$, $b' = 2\sin2\theta$, $c' =
\cos\theta$, dan $\Delta = 2\sin^22\theta + 2\cos^22\theta = 2$.
Lalu Cramer, kemudian rumus hasil-kali-ke-jumlah:
\begin{align*}
x &= \frac{2\sin\theta\sin2\theta +
\cos\theta\cos2\theta}{2}
= \frac{(\cos\theta - \cos3\theta) + \frac12(\cos\theta +
\cos3\theta)}{2} = \frac{3\cos\theta - \cos3\theta}4,\\
y &= \frac{\sin2\theta\cos\theta - 2\cos2\theta\sin\theta}2
= \frac{\frac12(\sin3\theta + \sin\theta) - (\sin3\theta -
\sin\theta)}2 = \frac{3\sin\theta - \sin3\theta}4 :
\end{align*}
yakni \omterm{pb:b2:curves:1}{nefroidnya}, sebuah kurva tertutup dengan dua titik runcing.

\textbf{19.} Menurunkannya lalu memfaktorkan dengan $\sin3\theta
- \sin\theta = 2\cos2\theta\sin\theta$ dan $\cos\theta -
\cos3\theta = 2\sin2\theta\sin\theta$:
\[
E'(\theta) = \tfrac34\bigl(\sin3\theta - \sin\theta,\
\cos\theta - \cos3\theta\bigr)
= \tfrac32\sin\theta\,(\cos2\theta, \sin2\theta),
\]
yang sejajar dengan arah pantul pada pertanyaan 16: jadi tiap
sinar pantulnya bersinggungan dengan kaustiknya, sebagaimana dituntut sifat
\omterm{pb:b2:curves:1}{selubungnya}. Adapun $E' = 0$ persis di $\sin\theta = 0$:
yakni $E(0) = (\tfrac12, 0)$ dan $E(\pi) = (-\tfrac12, 0)$, yaitu kedua
titik runcingnya. Lalu menetapkan $y = 0$ pada persamaan garisnya: $x\sin2\theta =
\sin\theta$, jadi $x = \frac1{2\cos\theta} \to \frac12$ ketika
$\theta \to 0$: sehingga sinar paraksialnya memfokus pada jarak $R/2$ dari
pusatnya --- yakni \omterm{def:b2:curves:length}{panjang} fokus cermin sferisnya.

\textbf{20.} Berlaku $\norm{E'(\theta)} = \tfrac32\abs{\sin\theta}$,
jadi \omterm{def:b2:curves:length}{panjangnya} $\tfrac32\int_0^{2\pi}\abs{\sin\theta}\,
\dd\theta = \tfrac32\cdot4 = 6$. Di dekat $\theta = 0$, dengan
$\cos k\theta = 1 - \tfrac{k^2\theta^2}2 + O(\theta^4)$ dan
$\sin k\theta = k\theta - \tfrac{k^3\theta^3}6 +
O(\theta^5)$:
\[
x - \tfrac12 = \frac{3\theta^2 + O(\theta^4)}{4}
= \tfrac34\theta^2 + O(\theta^4),
\qquad
y = \frac{4\theta^3 + O(\theta^5)}{4}
= \theta^3 + O(\theta^5) :
\]
jadi $p = 2$, $q = 3$, yakni titik runcing jenis pertama yang menunjuk sepanjang
sumbunya --- yaitu titik terang kaustik cangkir kopinya.

\textbf{21.} Parameterkan titik yang tersinari lewat
$\Phi(\theta, r) = P_\theta + r\,v_\theta$ (yakni kedudukan sepanjang
tiap sinar pantulnya). Adapun \omterm{def:b2:linalg:det}{determinan} Jacobinya
$\det(\partial_\theta\Phi, \partial_r\Phi) =
\det(P_\theta' + rv_\theta', v_\theta)$ bersifat \omterm{def:b2:affine:subspace}{afin} dalam $r$ dan
lenyap tepat untuk satu $r = r_*(\theta)$ --- lalu
$\Phi(\theta, r_*(\theta))$ adalah titik karakteristiknya, sebab
di sana arah sinarnya dan variasi keluarganya menjadi
bergantungan. Di luar \omterm{pb:b2:curves:1}{selubungnya} pemetaannya berupa difeomorfisma lokal,
dan titik tepat di dalam kaustiknya dikenai dua sinar berdekatan
(yakni dua penyelesaian $\theta$), sedangkan titik di luarnya tak dikenai satu pun dari
bagian keluarganya itu; sementara pada kaustiknya keduanya melebur. Adapun
intensitas cahayanya berbanding terbalik dengan nilai mutlak Jacobinya,
jadi ia meledak sepanjang \omterm{pb:b2:curves:1}{selubungnya}: sehingga kaustiknya adalah
kurva yang terang, dan paling terang di titik runcingnya, tempat
kemerosotannya paling parah.

\textbf{22.} Berlaku $x' = 1 - \cos t$, $y' = \sin t$, $x'' = \sin
t$, $y'' = \cos t$, jadi $x'y'' - y'x'' = \cos t - 1$ dan
$\norm{\gamma'}^2 = 2(1 - \cos t) = 4\sin^2\tfrac t2$;
sehingga menurut \cref{prop:b2:curves:kappaformula},
\[
\kappa(t) = \frac{-(1 - \cos t)}{8\sin^3\tfrac t2}
= \frac{-2\sin^2\tfrac t2}{8\sin^3\tfrac t2}
= -\frac1{4\sin\tfrac t2}
\qquad (0 < t < 2\pi).
\]
Dengan $T = (\sin\tfrac t2, \cos\tfrac t2)$ (yakni bagilah $\gamma'$
dengan $2\sin\tfrac t2$) dan $N = (-\cos\tfrac t2, \sin\tfrac
t2)$:
\[
\gamma + \frac1\kappa N
= \gamma - 4\sin\tfrac t2\,\Bigl(-\cos\tfrac t2,\
\sin\tfrac t2\Bigr)
= \bigl(t - \sin t + 2\sin t,\ 1 - \cos t - 2(1 - \cos
t)\bigr),
\]
yakni $c(t) = (t + \sin t,\ \cos t - 1)$. Setelah menyulihkan $t = u
+ \pi$:
\[
c = \bigl(u + \pi - \sin u,\ -\cos u - 1\bigr)
= \bigl((u - \sin u) + \pi,\ (1 - \cos u) - 2\bigr) :
\]
yakni sikloid $\gamma(u)$ yang digeser sebesar $(\pi, -2)$. Jadi
evolut sebuah sikloid adalah sikloid yang kongruen, yang tergantung satu
aras di bawahnya.

\textbf{23.} Berlaku $R(t) = 1/\abs{\kappa(t)} = 4\sin\tfrac t2$, jadi
$R(\pi) = 4$: yakni separuh \omterm{def:b2:curves:length}{panjang} lengkungnya $8$ yang dihitung pada
\cref{exo:b2:curves:1}. Adapun kecepatan evolutnya $c'(t) = (1 +
\cos t, -\sin t)$ lenyap di $t = \pi$: jadi titik runcingnya adalah $c(\pi)
= (\pi, -2)$, tepat di bawah puncaknya $\gamma(\pi) = (\pi,
2)$, pada jarak $4 = R(\pi)$, yakni persis \omterm{def:b2:curves:length}{panjang}
jari-jari oskulasinya di sana.

\textbf{24.} Dari \cref{exo:b2:curves:6}, berlaku $c'(s) =
-\frac{\kappa'(s)}{\kappa(s)^2}N(s) = R'(s)\,N(s)$, jadi
$\norm{c'(s)} = \abs{R'(s)}$ dan, untuk $R$ yang monoton,
\[
\int_{s_0}^{s_1}\norm{c'(s)}\,\dd s =
\Bigl|\int_{s_0}^{s_1}R'(s)\,\dd s\Bigr| = \abs{R(s_1) -
R(s_0)} .
\]
Katakanlah $R$ turun. Sebuah tali yang dibentangkan sepanjang evolutnya melewati
$c(s_0)$ lalu diperpanjang oleh ruas dari $c(s_0)$ ke
$\gamma(s_0)$ (yang bersinggungan dengan evolutnya, menurut pertanyaan 4)
punya, ketika dikupas sampai $c(s)$ lalu ditegangkan, bagian
lurus berpanjang $R(s_0) - \bigl(R(s_0) - R(s)\bigr) = R(s)$
yang menunjuk dari $c(s)$ sepanjang normalnya --- lalu mendarat tepat pada
$\gamma(s)$: jadi ujung bebasnya melacak kurva semula
(yakni ``involut''). Huygens menggantungkan sebuah bandul di antara dua pipi
sikloid: sehingga talinya membelit evolutnya, jadi bandulnya menggambarkan sebuah
sikloid --- yakni tautokron, yang kala ayunannya tak
bergantung pada amplitudonya.

\textbf{25.} Singgung parabolanya: $\Delta = 1$ dan $E' =
(1, t) \neq 0$ di mana-mana --- jadi \omterm{pb:b2:curves:1}{selubung} yang mulus (yakni parabolanya
sendiri). Tangga yang meluncur: $\Delta = -1$, tetapi $E' =
\tfrac32\sin 2t\,(-\cos t, \sin t)$ lenyap di ujung
kuadrannya --- yakni keempat titik runcing \omterm{pb:b2:curves:1}{astroidnya}. Normal
parabola dan sikloidnya: $\Delta = \kappa \neq 0$, dan
$E' = R'N$ lenyap persis di tempat \omterm{thm:b2:curves:frenet2d}{kelengkungannya} ekstrem
--- yakni titik runcing evolutnya di $(0,1)$ dan $(\pi, -2)$. Kaustiknya:
$\Delta = 2$, dan $E' = \tfrac32\sin\theta\,(\cos2\theta,
\sin2\theta)$ lenyap di $\theta = 0, \pi$ --- yakni kedua titik runcing
\omterm{pb:b2:curves:1}{nefroidnya}, yaitu titik fokus cerminnya. Adapun ekstremum
\omterm{thm:b2:curves:frenet2d}{kelengkungan} \emph{mesti} menghasilkan titik runcing pada \omterm{pb:b2:curves:1}{selubung}
normalnya, karena kecepatan evolutnya adalah $R'N$: dan itulah sebabnya
evolut sebuah elips punya empat titik runcing (yakni empat puncak),
sedangkan keluarga yang merosot (berkas, garis sejajar) adalah kasus
tempat mesinnya mengeluarkan sebuah titik atau tak ada apa pun.

\end{solution}
